\chapter{Diskussion}

Die Diskussion ist das intellektuell anspruchsvollste Kapitel. Hier werden die Ergebnisse interpretiert, mit der Literatur verglichen und kritisch eingeordnet. Während das Ergebniskapitel fragt \glqq Was kam heraus?\grqq, fragt die Diskussion \glqq Was bedeutet das?\grqq.

\section{Interpretation der Resultate}

Wurden die in der Einleitung definierten Ziele erreicht? Wie verhalten sich die eigenen Ergebnisse im Vergleich zu den im Stand der Technik beschriebenen Arbeiten? Wenn es Abweichungen zwischen Erwartung und Realität gab — woran könnten diese liegen?

Hier dürfen und sollen Hypothesen aufgestellt werden, solange sie als solche gekennzeichnet sind. Formulierungen wie \glqq Eine mögliche Erklärung ist\ldots\grqq{} oder \glqq Dies könnte darauf hindeuten, dass\ldots\grqq{} sind typisch für wissenschaftliches Schreiben.

\section{Kritische Würdigung}

Keine Arbeit ist perfekt — und das offen zu benennen ist eine Stärke, keine Schwäche. Typische Fragen sind: Welche Annahmen mussten getroffen werden? Wie groß war die Stichprobe? Welche Fehlerquellen gibt es? Was würde man beim nächsten Mal anders machen?

Eine ehrliche kritische Würdigung zeigt dem Prüfer, dass der Verfasser die Grenzen der eigenen Arbeit versteht und wissenschaftlich reflektiert.